\section*{第 \ref{ch:g12:sums} 章 --- 随机变量之和与大数定律}

\begin{solution}{exo:g12:sums:1}
每个骰子的期望为
$\frac{1 + 2 + \dots + 6}{6} = \frac{7}{2}$，由线性性
可得 $\E(S) = \frac72 + \frac72 = 7$。骰子相互独立，于是
$\V(S) = \frac{35}{12} + \frac{35}{12} = \frac{35}{6} \approx 5.83$。
\end{solution}

\begin{solution}{exo:g12:sums:2}
$\E(X) = 50$ 且 $\V(X) = 100 \times 0.25 = 25$。

马尔可夫不等式（$X \geq 0$）：$\P(X \geq 75) \leq \frac{50}{75} = \frac23$。

比内梅--切比雪夫不等式：
$\P(\abs{X - 50} \geq 25) \leq \frac{25}{25^2} = 0.04$ --- 这一估计甚至
控制的是\emph{双侧}事件，其中 $X \geq 75$ 仅占一半。切比雪夫
在此精确得多，因为它利用了方差，而不只是均值。
（$\P(X \geq 75)$ 的真实值小于 $10^{-6}$：两个界都
很粗糙。）
\end{solution}

\begin{solution}{exo:g12:sums:3}
$F_n$ 是 $n$ 个伯努利 $\left(\frac12\right)$ 变量的样本均值，
其方差为 $\sigma^2 = \frac14$。下式
\[
\P\left(\abs{F_n - 0.5} \geq 0.05\right)
\leq \frac{1/4}{n \times 0.05^2} = \frac{100}{n}
\]
该式 $\leq 0.05$ 只需 $n \geq 2000$。
\end{solution}

\begin{solution}{exo:g12:sums:4}
由\cref{prop:g12:sums:mean}可知，$\E(M_n) = 5\%$（分散化并不
改变期望收益），且
$\sigma(M_n) = \frac{20\%}{\sqrt n}$。要求
$\frac{20}{\sqrt n} < 4$ 得 $\sqrt n > 5$，\emph{即} $n \geq 26$。
这就是\emph{分散化}原理：将资金分散到
相互独立的风险上，使风险（标准差）除以 $\sqrt n$，
而不降低期望收益。
\end{solution}

\begin{solution}{exo:g12:sums:5}
\emph{1.} 数出 $9$ 个等可能的数对：$S$ 取值
$2, 3, 4, 5, 6$ 的概率分别为
$\frac19, \frac29, \frac39, \frac29, \frac19$。于是
$\E(S) = \frac{2 + 6 + 12 + 10 + 6}{9} = 4$，且
$\E(S^2) = \frac{4 + 18 + 48 + 50 + 36}{9} = \frac{156}{9} = \frac{52}{3}$，
从而 $\V(S) = \frac{52}{3} - 16 = \frac43$。

\emph{2.} 考虑一个变量：$\E(X) = 2$，
$\E(X^2) = \frac{1 + 4 + 9}{3} = \frac{14}{3}$,
$\V(X) = \frac{14}{3} - 4 = \frac23$。于是 $\E(S) = 2 + 2 = 4$，再由
独立性，$\V(S) = \frac23 + \frac23 = \frac43$。两者的值相同。
\end{solution}

\begin{solution}{exo:g12:sums:6}
取 $Y = X$，则 $\V(X + Y) = \V(2X) = 4\V(X)$，而
$\V(X) + \V(Y) = 2\V(X)$。二者仅当 $\V(X) = 0$ 时相等，\emph{即}
当 $X$ 为常数时 --- 因此可加性确实需要独立性
（此处 $X$ 与其自身完全相依）。
\end{solution}

\begin{solution}{exo:g12:sums:7}
在公平假设下，$F_n$ 是 $n = 1200$ 个伯努利变量的均值，其中
$p = \frac16$，$\sigma^2 = \frac16\cdot\frac56 = \frac{5}{36}$：
\[
\P\left(\abs{F_n - \tfrac16} \geq 0.05\right)
\leq \frac{5/36}{1200 \times 0.0025} = \frac{5}{108} \approx 0.046 .
\]
观测偏差恰为 $0.05$：这一事件由公平骰子产生
的概率至多为 $4.6\%$ --- 由于切比雪夫不等式十分保守，
真实概率远小于此。公平性假设
并不可信，这颗骰子很可能是有偏的。
\end{solution}

\begin{solution}{exo:g12:sums:8}
\emph{1.} $p(1-p) = \frac14 - \left(p - \frac12\right)^2 \leq \frac14$。

\emph{2.} $F_n$ 的期望为 $p$，方差为
$\frac{p(1-p)}{n} \leq \frac{1}{4n}$；由比内梅--切比雪夫不等式得
\[
\P\bigl(\abs{F_n - p} \geq \varepsilon\bigr)
\leq \frac{p(1-p)}{n\varepsilon^2} \leq \frac{1}{4n\varepsilon^2},
\]
对任意未知的 $p$ 均成立。

\emph{3.} 取 $\varepsilon = 0.03$，误差水平为 $0.05$：
\[
\frac{1}{4n(0.03)^2} \leq 0.05
\iff n \geq \frac{1}{4 \times 0.0009 \times 0.05} \approx 5556 .
\]
按此粗略界，约 $5600$ 人即足够（经典的
正态近似结果更接近 $1100$，见
\cref{ch:g12:contdist}）。
\end{solution}

\begin{solution}{pb:g12:sums:1}
\textbf{1.} $V(X + Y) = V(X) + V(Y) = \frac{35}{12} +
\frac{35}{12} = \frac{35}{6}$（由独立性）。

\textbf{2.} $\E(\bar X) = 3.5$;
$\sigma(\bar X) = \frac{1.71}{\sqrt{100}} \approx 0.17$：百颗骰子的
平均值紧贴 $3.5$，偏差在五分之一分以内。

\textbf{3.} $\P(X \geq 10) \leq \frac{2}{10} = 0.2$.

\textbf{4.} $\P \leq \frac{V(\bar X)}{0.5^2} =
\frac{35/1200}{0.25} \approx 0.117$：至多约
$12\,\%$。

\textbf{5.} 令 $Y = -X$：$V(X + Y) = V(0) = 0$，而
$V(X) + V(Y) = 2V(X) > 0$：完全负相关的变量
相互抵消 --- 方差的可加性是独立性赋予的特权。

\textbf{6.} $\E = \frac{18}{37} - \frac{19}{37} =
-\frac{1}{37} \approx -0.027$;
$\sigma = \sqrt{1 - \left(\frac{1}{37}\right)^2} \approx
1.00$.

\textbf{7.} $\E(G) = -2.70$, $\sigma(G) = 10 \times 1.00 =
10$：漂移仅为 $0.27\sigma$。单晚的噪声远远淹没
优势 --- 相当多的赌徒（约 $40\,\%$，
钟形曲线所言）赢着离场，
这恰恰让他们一再回头。

\textbf{8.} 赌场的收益：期望 $+27\,027$ 欧元，
标准差 $\approx 1\,000$：利润位于零之上 $27$
个标准差。在 $27\sigma$ 处，“赌场
今年可能亏损”就不是风险，而是舍入误差。

\textbf{9.} $\P(\text{亏损}) = \P(G \geq 27\,027$ 对赌徒
而言$)$ $\leq \frac{10^6 \times 1}{27\,027^2} \approx
0.0014$：即便是书中最粗糙的不等式，也保证了
庄家有 $99.86\,\%$ 的胜算 --- 真实情况
要强得惊人。

\textbf{10.} 每投入一欧元，无论何时、以何种方式选择，
其期望回报为自身的 $-\frac{1}{37}$；由期望的线性性，
总期望收益为 $-\frac{1}{37} \times
(\text{总投注额})$，对任何下过注的策略都为负。
一个能获利的系统，将违背期望的线性性
--- 再巧妙地排列坏赌注的次序，也造不出
一个好的。（加倍投注的系统，不过是以许多小赢
换取少见的灾难性亏损。）

\textbf{11.} 定理断言：红色的 \emph{频率}
在 $n$ 次转动中依概率收敛于 $\frac{18}{37}$。它
对第 $n + 1$ 次转动没有任何断言：转盘没有记忆，而
即便已连续十次出红，红色的机会仍是 $\frac{18}{37}$。
持相反看法即犯了\emph{赌徒谬误}，而
第 12 题说明了连续同色现象的真实走向。

\textbf{12.} $100$ 次红色的盈余不会被偿还 ---
未来的转动都公平独立，期望盈余仍为 $100$。
但 $\frac{100}{10^6} = 0.0001$：仅为一个百分点的百分之一。
该谬误的错误可用一句话概括：\emph{大数定律
以新试验的汪洋稀释过去的偶然；它
从不会让转盘去讨还债务。}

\textbf{13.} $n \geq \frac{1}{4 \times 0.05 \times 0.03^2}
\approx 5\,556$ 人。

\textbf{14.} $\frac{1.96^2 \times 0.25}{0.03^2} \approx
1\,067$：约少了五倍。切比雪夫不等式对\emph{任何}
分布都成立，为其普适性付出了松弛的代价；
而民调常数 $1.96$ 来自涨落的实际钟形形状，
后者的集中程度远为紧密。

\textbf{15.} 误差幅度 $\propto \frac{1}{\sqrt n}$：将其减半
需要样本量增至四倍 --- 约 $4\,400$ 人才能给出
$\pm 1.5$ 个百分点的精度。精度是以平方级的代价换来的。

\textbf{16.} 模型是 $n$ 次独立抽样，每次成功的
概率为 $p$ --- 总体规模 $N$ 从不出现
（从一个大而充分混合的总体中抽取极小的一部分，
正是独立性所刻画的内容）。民调者的比方是：
要尝汤，搅匀后的一勺便已足够 ---
无论锅里盛的是十人份还是一万人份。关键在于
搅拌 --- 抽样框有偏差，正如初中分册中
那场骗过了全国的民调，而不在锅的大小。

\textbf{17.} $315\sqrt n < 30n \iff \sqrt n > 10.5 \iff
n > 110$：超过一百份保单后，漂移便压倒
波动，每多一份保单，差距便按
$\sqrt n$ 律扩大 --- 汇集风险\emph{即是}保险的商业模式。

\textbf{18.} 独立等额持仓：$\sigma_{\text{组合}} =
\frac{20\,\%}{\sqrt{25}} = 4\,\%$：期望收益相同，
波动仅为五分之一。这顿“免费的午餐”：风险降低，
期望上没有任何代价 --- 即对独立的随机性取平均。

\textbf{19.} 完全正相关：该组合实为同一只股票
穿上 $25$ 套“行头”：$\sigma = 20\,\%$，毫无降低。
分散化租借的是\emph{独立性}；一旦危机
使一切都相关（2008 年：所有房产类赌注实为同一赌注），
$\sqrt n$ 的保护恰在
最需要的时候消失。

\textbf{20.} 和的漂移按 $n$ 增长，波动按
$\sqrt n$ 增长，故均值按 $\frac{1}{\sqrt n}$ 的速率稳定下来：
赌场在百万次转动中收取漂移；民调者
以电话预算换得 $\frac{1}{\sqrt{1100}}$ 的噪声；
保险公司在 $n > 110$ 时便超过了自身的理赔；
投资者将风险除以 $\sqrt{25}$。误用者包括：赌徒，
相信债务必须偿还（其实只有稀释），以及
相信独立性的金融家（有时实际上
只有一个赌注）。涨落的普遍形状 ---
钟形曲线 --- 是连续型那一章的主角，而其定理
即中心极限定理，是本科分册中
概率论的巅峰。
\end{solution}
